\lesson{7}{от 3 июня 2025 12:45}{Продолжение}

\begin{definition}[Нейронная сеть]
  \emph{Нейронная сеть} -- математическая модель, а также ее программное или аппаратное воплощение, построенная по принципу организации биологических нейронных сетей -- сетей нервных клеток живого организма.
\end{definition}

\begin{note}
  После разработки алгоритмов обучения получаемые модели стали использовать в практических целях: в задачах прогнозированния, для распознавания образов, в задачах управления и т.д.
\end{note}

\begin{enumerate}
  \item С точки зрения машинного обучения, нейронная сеть представляет собой частный случай методов распознования образов.
    \item С точки зрения математики, обучение нейронных сетей -- это многопараметрическая задача нелинейной оптимизации.
\end{enumerate}

\begin{eg}
  \emph{Простой пример} нейронной сети, которая на входе получает рост и вес человека, а на выходе предсказывает пол. Скрытый слой в данном случае мы использовать не будем. В качестве функции активации возьмем сигмоиду. Ее часто используют в задачах бинарной (состоящей из двух классов) классификации.
\end{eg}

\section{Обучение нейронной сети}

Как и с традиционными алгоритмами ML, для построения модели, необходимо подобрать оптимальные веса или заняться оптимизацией. Применительно к нейронным сетям этот процесс называется \emph{обратным распросранением}.

\begin{note}
  Мы как бы двигаемся в обратную сторону и, уже зная результат (и уровень ошибки), с учетом имеющихся данных, рассчитываем, как нам нужно изменить веса и смещения, чтобы уровень ошибки снизился.
\end{note}

Если \emph{сигмоида} подходит для бинарной классификации, то \emph{softmax} применяется для задач многоклассовой классификации.

На выходе получается вероятностное распределение любого количества классов (K).

\begin{note}
  Настроек три:

  \begin{itemize}
    \item тип функции определяет, как мы будем считать отклонение прогнозного значения от истинного;
      \item способ или алгоритм оптимизации этой функции поможет снизить потерю или ошибку и подобрать правиьные веса в процессе back propagation;
        \item метрика покажет, насколько точна созданная модель.
  \end{itemize}
\end{note}

\section{Компьютерное зрение}

\begin{definition}[Компьютерное зрение]
  \emph{Компьютерное зрение} -- это область, которая занимаеся анализом фото или видео изображений. Построенные на основе этого анализа модели помогают смартфонам узнавать своих владельцев, дорожным камерам распознавать номера автомобилей, а роботам объезжать препятствия.
\end{definition}

\begin{note}
  Любое цифровое изображение -- это по сути матрица, а в случае Питона, массив Numpy, в котором \emph{пикселей} (составных частей) задано...
\end{note}

В случае черно-белого изображения, или правильнее сказать, изображения в оттенках серого, речь идет об одном слое матрицы.

\section{Нейросети и их типы}

\begin{itemize}
  \item Перцептрон: Основа Искусственного Интелекта;
    \item Многослойные Перцептроны (MLP): Строим Будущее С Искусственным Интелектом;
      \item Рекуррентные Нейронные Сети (RNN): Магия Последовательностей (временные ряды, тексты);
        \item Сверточные Нейронные Сети (CNN): Глаза Искусственного Интелекта (компьютерное зрение).
\end{itemize}
